\section{Fazit}

Ziel dieser Arbeit war die Übertragung der interaktiven Python-basierten Lehrnotebooks der Vorlesung „Theoretische Informatik III“ in das TypeScript-Ökosystem. Dabei wurde bewusst nicht nur eine syntaktische Übersetzung angestrebt, sondern eine strukturelle Neuausrichtung der zugrunde liegenden Architektur. Im Zentrum stand die Frage, inwiefern sich Konzepte formaler Sprachen durch statische Typisierung präziser modellieren und zuverlässiger umsetzen lassen.

Die Ergebnisse der Arbeit zeigen deutlich, dass TypeScript durch sein statisches Typsystem einen wesentlichen Mehrwert für die Modellierung komplexer Datenstrukturen bietet. Während Python viele Fehler erst zur Laufzeit sichtbar macht, ermöglicht TypeScript eine frühzeitige Validierung von Invarianten bereits während der Entwicklung. Insbesondere bei verschachtelten Strukturen wie Syntaxbäumen oder Zustandsmengen trägt dies maßgeblich zur Robustheit der Implementierung bei. Damit wird ein zentraler Vorteil statischer Typisierung bestätigt, nämlich die Vorverlagerung von Fehlererkennung in die Entwicklungsphase \autocite[S.~371ff]{aho2006compilers}.

Ein zentrales Ergebnis der Arbeit ist die Erkenntnis, dass eine naive Portierung nicht ausreicht. Unterschiede in der Laufzeitsemantik, wie etwa das referenzbasierte Verhalten von Mengen in JavaScript, erfordern eigenständige architektonische Lösungen. Die entwickelte Datenstruktur \texttt{RecursiveSet} zeigt exemplarisch, wie durch gezielte Implementierungsentscheidungen mathematische Konzepte wie Wertgleichheit und Immutabilität in TypeScript nachgebildet werden können. Gleichzeitig verdeutlicht sie, dass Typsysteme allein nicht alle semantischen Unterschiede zwischen Programmiersprachen überbrücken können, sondern durch geeignete Datenstrukturen ergänzt werden müssen.

Auch im Bereich der Sprachverarbeitung konnte gezeigt werden, dass die grundlegenden Prinzipien der Compilertechnik unabhängig von der Zielplattform bestehen bleiben. Die Transformation von der lexikalischen Analyse über die syntaktische Analyse bis hin zur semantischen Interpretation folgt sowohl in der Python- als auch in der TypeScript-Implementierung denselben theoretischen Grundlagen \autocite[S.~4ff]{aho2006compilers}. Der Einsatz von Lezer als Parsergenerator verdeutlicht jedoch, wie sich diese klassischen Konzepte an moderne Anforderungen wie inkrementelles Parsing und editornahe Anwendungen anpassen lassen.

Die Case Study zur Implementierung eines vollständigen Interpreters zeigt schließlich, dass sich die gesamte Verarbeitungskette – von der Grammatikdefinition bis zur Ausführung – erfolgreich in TypeScript abbilden lässt. Dabei wird deutlich, dass die statische Typisierung nicht nur der Fehlervermeidung dient, sondern auch die Struktur und Lesbarkeit des Codes verbessert. Gleichzeitig bleibt ein gewisser Trade-off bestehen: Eine vollständig präzise Typmodellierung kann zu erhöhter Komplexität führen und erfordert bewusste Entwurfsentscheidungen zwischen Strenge und Flexibilität.

Zusammenfassend lässt sich festhalten, dass die Übersetzung der Lehrmaterialien nach TypeScript nicht nur technisch möglich, sondern auch didaktisch sinnvoll ist. Die stärkere Typisierung unterstützt das Verständnis formaler Strukturen, zwingt zur expliziten Modellierung von Invarianten und reduziert typische Fehlerquellen dynamischer Systeme. Gleichzeitig zeigt die Arbeit, dass eine erfolgreiche Migration stets eine konzeptionelle Anpassung erfordert und nicht auf eine reine Syntaxübertragung reduziert werden kann.

Die entwickelten Ansätze bilden damit eine solide Grundlage für zukünftige Erweiterungen, etwa die weitergehende Typisierung von AST-Strukturen oder die Integration zusätzlicher Analyse- und Optimierungsschritte innerhalb der Interpreterpipeline.